% CUMCM (国赛) LaTeX Template
% Math Modeling Contest Agent
\documentclass[12pt,a4paper]{ctexart}

% === Packages ===
\usepackage{amsmath,amssymb,amsfonts}
\usepackage{graphicx}
\usepackage{booktabs}
\usepackage{multirow}
\usepackage{longtable}
\usepackage{enumitem}
\usepackage{hyperref}
\usepackage{geometry}
\usepackage{fancyhdr}
\usepackage{float}
\usepackage{caption}
\usepackage{subcaption}
\usepackage{algorithm}
\usepackage{algpseudocode}
\usepackage{xcolor}
\usepackage{listings}

% === Page Setup ===
\geometry{left=2.5cm,right=2.5cm,top=2.5cm,bottom=2.5cm}
\setlength{\parskip}{0.5em}
\setlength{\parindent}{2em}

% === Header/Footer ===
\pagestyle{fancy}
\fancyhf{}
\fancyhead[C]{\small 数学建模竞赛论文}
\fancyfoot[C]{\thepage}
\renewcommand{\headrulewidth}{0.4pt}

% === Code Style ===
\lstset{
    basicstyle=\ttfamily\small,
    breaklines=true,
    frame=single,
    numbers=left,
    numberstyle=\tiny,
    xleftmargin=2em,
}

% === Title ===
\title{\textbf{题目：\hspace{2em}\makebox[8cm]{\dotfill}}}
\author{}
\date{}

\begin{document}

\maketitle
\thispagestyle{empty}

% === Abstract ===
\begin{center}
\textbf{\Large 摘\quad 要}
\end{center}

% TODO: 摘要由 Agent 自动生成
本文针对[题目背景]，建立了[模型类型]模型，解决了[核心问题]。

\textbf{关键词：}关键词1；关键词2；关键词3；关键词4

\newpage
\tableofcontents
\newpage

% ==========================================
% 1. 问题重述
% ==========================================
\section{问题重述}

% TODO: Agent 自动生成

% ==========================================
% 2. 问题分析
% ==========================================
\section{问题分析}

% TODO: Agent 自动生成

% ==========================================
% 3. 模型假设与符号说明
% ==========================================
\section{模型假设与符号说明}

\subsection{基本假设}
\begin{enumerate}[label=(\arabic*)]
    \item 假设1
    \item 假设2
\end{enumerate}

\subsection{符号说明}
\begin{table}[H]
\centering
\caption{符号说明表}
\begin{tabular}{cll}
\toprule
\textbf{符号} & \textbf{含义} & \textbf{单位} \\
\midrule
$x$ & 变量说明 & - \\
\bottomrule
\end{tabular}
\end{table}

% ==========================================
% 4. 模型建立与求解
% ==========================================
\section{模型建立与求解}

\subsection{问题一：模型名称}
\subsubsection{模型建立}
% TODO: Agent 自动生成数学模型推导

\subsubsection{模型求解}
% TODO: Agent 自动生成求解过程和结果

\subsubsection{结果分析}
% TODO: Agent 自动生成结果分析

\subsection{问题二：模型名称}
% 结构同上

\subsection{问题三：模型名称}
% 结构同上

% ==========================================
% 5. 灵敏度分析
% ==========================================
\section{灵敏度分析}
% TODO: Agent 自动生成

% ==========================================
% 6. 模型评价与改进
% ==========================================
\section{模型评价与改进}

\subsection{模型优点}
\begin{enumerate}
    \item 优点1
    \item 优点2
\end{enumerate}

\subsection{模型缺点}
\begin{enumerate}
    \item 缺点1
    \item 缺点2
\end{enumerate}

\subsection{改进方向}
% TODO: Agent 自动生成

% ==========================================
% 参考文献
% ==========================================
\begin{thebibliography}{99}

\bibitem{ref1} 作者. 标题[J]. 期刊, 年份, 卷(期): 页码.
\bibitem{ref2} 作者. 书名[M]. 出版地: 出版社, 年份.

\end{thebibliography}

% ==========================================
% 附录
% ==========================================
\appendix
\section{核心代码}
% 附上关键求解代码

\end{document}
